(1) 理想弦满足 $\mu y_{tt}=Ty_{xx}$，波速 $c=\sqrt{T/\mu}$。d'Alembert公式 $y(x,t)=\frac12[y_0(x-ct)+y_0(x+ct)]+(2c)^{-1}\int_{x-ct}^{x+ct}v_0(s)ds$，$v_0=0$，所以 $y(x,t)=\frac A2[e^{-(x-ct)^2/(2a^2)}+e^{-(x+ct)^2/(2a^2)}]$。

(2) 两脉冲峰值在 $x=\pm ct$，分离后的峰高各近 $A/2$，原点 $y(0,t)=A e^{-c^2t^2/(2a^2)}$。每个脉冲的函数形状始终只发生平移，没有改变高斯宽度，因为色散律为 $\omega=c|k|$，没有群速度色散。

(3) 乘波动方程以 $y_t$ 得 $\partial_t[\mu y_t^2/2+Ty_x^2/2]+\partial_x[-Ty_xy_t]=0$。初始动能零，$E_0=\frac T2\int y_0'^2dx=\frac{TA^2}{2a^4}\int x^2e^{-x^2/a^2}dx=TA^2\sqrt\pi/(4a)$。

(4) 足够久之后两脉冲空间分离，每个单向脉冲 $y_\pm=(A/2)e^{-(x\mp ct)^2/(2a^2)}$，动能密度和张力势能密度相等，所以其总能量为 $E_\pm=T\int (\partial_xy_\pm)^2dx=TA^2\sqrt\pi/(8a)$；两份相加 $E_++E_-=E_0$。任意有限时两份场的交叉项仍要合并算总能量，不能把尚未分离的脉冲各自视为完全独立。
